% Source: Weinberg GR, physical PDF pp. 8--25, plus supplemental Contents
%         p. xxii (`source-supplements/contents-pxxii.png`) between physical
%         pp. 19--20
%         (printed p. ix and Contents pp. xi--xxvii).
% Coverage: Modernized notation leaf and source contents; master.tex generates
%           the typeset table of contents from the complete assembly.
% Figures/tables/footnotes: none.
% Status: source-reviewed and compile-clean.
% Uncertainties: none.

\section*{Notation}
\addcontentsline{toc}{section}{Notation}

This edition modernizes the notation of the original text according to the
following conventions.

Greek indices $\mu,\nu,\rho,\sigma,\ldots$ run over the four space-time
coordinates $0,1,2,3$, with
\[
  x^\mu=(x^0,x^1,x^2,x^3)=(t,\mathbf{x}).
\]
Latin indices $i,j,k,l,\ldots$ run over the three spatial coordinates
$1,2,3$. Latin indices $a,b,c,\ldots$ denote local Lorentz indices when a
tetrad is used. Repeated indices are summed unless otherwise indicated.

The Minkowski metric has mostly-plus signature,
\[
  \eta_{\mu\nu}=\diag(-1,+1,+1,+1),
\]
and a timelike four-velocity obeys $u^\mu u_\mu=-1$. Our curvature convention
is
\[
  [\nabla_\mu,\nabla_\nu]V^\rho
  =R^\rho{}_{\sigma\mu\nu}V^\sigma,
\]
\[
  R^\rho{}_{\sigma\mu\nu}
  =\partial_\mu\Gamma^\rho{}_{\nu\sigma}
  -\partial_\nu\Gamma^\rho{}_{\mu\sigma}
  +\Gamma^\rho{}_{\mu\lambda}\Gamma^\lambda{}_{\nu\sigma}
  -\Gamma^\rho{}_{\nu\lambda}\Gamma^\lambda{}_{\mu\sigma}.
\]
Thus
\[
  G_{\mu\nu}+\Lambda g_{\mu\nu}=8\pi G\,T_{\mu\nu}.
\]

Partial and covariant derivatives are denoted by $\partial_\mu$ and
$\nabla_\mu$. Weinberg's notation for covariant differentiation along a
worldline is retained:
\[
  \frac{D V^\mu}{D\tau}\equiv u^\nu\nabla_\nu V^\mu.
\]
A dot denotes differentiation with respect to time. Cartesian three-vectors
are indicated by boldface type.

The numerical permutation symbol is denoted by
$\tilde\epsilon^{0123}=+1$, while the Levi--Civita tensor is denoted by
\[
  \varepsilon^{\mu\nu\rho\sigma}
  =\frac{\tilde\epsilon^{\mu\nu\rho\sigma}}{\sqrt{-g}},
  \qquad
  \varepsilon_{\mu\nu\rho\sigma}
  =-\sqrt{-g}\,\tilde\epsilon_{\mu\nu\rho\sigma}.
\]
Thus $\varepsilon_{0123}=-\sqrt{-g}$ in the mostly-plus convention. A fixed
orientation is understood; orientation-reversing maps give the usual
pseudotensor sign.

The action is denoted by $I$, as in the original:
\[
  I=I_M+I_G,\qquad
  I_G=\frac{1}{16\pi G}\int d^4x\,\sqrt{-g}\,(R-2\Lambda).
\]
The speed of light is taken to be unity except when c.g.s.\ units are
indicated. Planck's constant is not taken to be unity.
